\section{Introduction and main results}
\label{sec:introduction}

Point interactions on a chain of shrinking intervals can remain visible
in the second term of a high-energy counting law. On the harmonic chain,
the critical balance occurs when the interaction at the $n$th vertex
has strength comparable to $n$. We prove that the resulting second
coefficient is determined by the full periodic-cell integrated density
of states, rather than by its free or Dirichlet endpoint.

The harmonic-chain model is classical. Egger n\'e Endres and
Steiner~\cite{egger2011trace} considered constant positive coupling,
proved discreteness, and analyzed in detail the fully decoupled
Dirichlet limit. Its eigenvalues are the squares of positive integers
with divisor multiplicities, so its counting law is an instance of
Dirichlet's divisor problem. General realization and discreteness
results for shrinking-gap point interactions are developed by Kostenko
and Malamud~\cite{kostenko2010point}. We include a direct form argument
because its quantitative tail bound is needed for the counting proof,
not because compactness itself is new.

The periodic comparison uses the usual delta-jump and free-propagation
matrices; their constant-coupling specialization and band structure
are recalled, for example, in Drabkin, Kirsch, and
Schulz-Baldes~\cite[Secs.~2.1--2.3]{drabkin2012transport}. The spectral
comparison work of Bifulco and Kerner~\cite{bifulco2024comparison}
also revisits the harmonic chain. Its modified local Weyl theorem
in \emph{arXiv v1, Section~5, Theorem~11} assumes fully Dirichlet
coupling. That local theorem does not provide the global critical
second coefficient considered here. These comparisons identify the
background used below; we make no general priority claim concerning
all slowly varying periodic-operator asymptotics.

\subsection{The source operator and the periodic comparison}

Write
\begin{equation}
 x_0=0,\qquad x_n=\pi\harm_n,
 \qquad \harm_n=\sum_{j=1}^n\frac1j,
 \qquad I_n=(x_{n-1},x_n).
 \label{eq:geometry}
\end{equation}
Thus $\abs{I_n}=\pi/n$, while $x_n\to\infty$.
Throughout the paper, $(b_n)_{n\ge1}$ is a sequence of positive real
numbers, fixed independently of the spectral frequency $k$, and
\begin{equation}
 n b_n\longrightarrow\infty.
 \label{eq:compact-assumption}
\end{equation}
Let
\begin{equation}
 \begin{split}
 \V&=\{f\in H^1(0,\infty):f(0)=0\},\\
 q_b[f]&=\int_0^\infty\abs{f'(x)}^2\dd x
            +\sum_{n\ge1}b_n\abs{f(x_n)}^2,\\
 \D_b&=\left\{f\in\V:
                 \sum_{n\ge1}b_n\abs{f(x_n)}^2<\infty\right\}.
 \end{split}
 \label{eq:source-form}
\end{equation}
Only the trace at zero is prescribed; there is no additional boundary
condition at infinity. Section~\ref{sec:forms} proves that this form is
densely defined and closed, with compact embedding into $L^2(0,\infty)$.
Its represented self-adjoint operator is denoted by $H_b$.
It is the form realization of
$-\mathrm d^2/\mathrm dx^2+\sum_n b_n\delta_{x_n}$, with continuity
and derivative jump $f'(x_n+)-f'(x_n-)=b_n f(x_n)$ at each vertex.
The associated source dynamics is $\exp(-\ii tH_b)$; our observable is
the inclusive spectral count
\begin{equation}
 N_b(k^2)=\#\{j:\lambda_j(H_b)\le k^2\},\qquad k\ge0,
 \label{eq:inclusive-count}
\end{equation}
where eigenvalues are counted with multiplicity.

For $\kappa>0$ and $\theta\in[0,2\pi)$, define the comparison-cell form
\begin{equation}
 \begin{split}
 \D_\theta&=\{u\in H^1(0,\pi):
                       u(\pi)=e^{\ii\theta}u(0)\},\\
 q_{\kappa,\theta}[u]
    &=\int_0^\pi\abs{u'(x)}^2\dd x+\kappa\abs{u(0)}^2.
 \end{split}
 \label{eq:fiber-form}
\end{equation}
Let $\nu_\kappa(z,\theta)$ count its eigenvalues at most $z^2$, and put
\begin{equation}
 J_\kappa(z)=\frac1{2\pi}\int_0^{2\pi}
                   \nu_\kappa(z,\theta)\dd\theta,
 \qquad z\ge0.
 \label{eq:ids-definition}
\end{equation}
This is the integrated density of states \emph{per cell of length
$\pi$}, not per unit length. The periodic operator is a local comparison
problem, not a claimed global conjugate of $H_b$.

\begin{theorem}[Critical second coefficient]
\label{thm:critical}
Under \eqref{eq:compact-assumption}, $H_b$ has strictly positive
discrete spectrum. If
\begin{equation}
 \frac{b_n}{n}\longrightarrow\kappa\in(0,\infty),
 \label{eq:critical-assumption}
\end{equation}
then, as real $k\to\infty$,
\begin{equation}
 N_b(k^2)=k\log k+C(\kappa)k+o(k),
 \label{eq:main-law}
\end{equation}
where
\begin{equation}
 C(\kappa)=\gamma+
     \int_0^\infty
       \frac{J_\kappa(z)-z\ind_{[1,\infty)}(z)}{z^2}\dd z.
 \label{eq:coefficient-definition}
\end{equation}
The integral converges. The function $C:(0,\infty)\to\R$ is
continuous and strictly decreasing, with
\begin{equation}
 C((0,\infty))=(2\gamma-1,\infty),\qquad
 \lim_{\kappa\to\infty}C(\kappa)=2\gamma-1,
 \qquad \lim_{\kappa\downarrow0}C(\kappa)=\infty.
 \label{eq:coefficient-range}
\end{equation}
No rate or monotonicity is required in \eqref{eq:critical-assumption}.
In particular, the conclusion holds for every positive sequence
$b_n=\kappa n+o(n)$.
\end{theorem}

\begin{theorem}[Hard and soft limits]
\label{thm:endpoint}
Retain positivity of $(b_n)$ and \eqref{eq:compact-assumption}.
If $b_n/n\to\infty$, then
\begin{equation}
 N_b(k^2)=k\log k+(2\gamma-1)k+o(k).
 \label{eq:hard-law}
\end{equation}
If $b_n/n\to0$, then the conclusion is precisely
\begin{equation}
 \frac{N_b(k^2)-k\log k}{k}\longrightarrow\infty.
 \label{eq:soft-law}
\end{equation}
All limits include eigenvalue thresholds. In the second case no
general leading-order formula for $N_b(k^2)$ is asserted.
\end{theorem}

Fixed positive couplings $b_n\equiv b$ belong to the soft side,
not to \eqref{eq:critical-assumption}. Passing to a finite nonzero
limit of $b_n/n$ changes the asymptotic balance: for indices $n$ of
order $k$, the product $b_n\abs{I_n}$ approaches $\pi\kappa$.
Finite transmission therefore persists on an order-$k$ set of cells.
An endpoint replacement on those cells can change the coefficient
of $k$ even when it leaves the leading term unchanged.

\subsection{The counting mechanism}

The principal analytic assertion is
\begin{equation}
 N_{(\kappa n)}(k^2)=\sum_{n\ge1}J_\kappa(k/n)+o(k).
 \label{eq:intro-reduction}
\end{equation}
Here $N_{(\kappa n)}$ denotes the count for $b_n=\kappa n$ exactly.
A bounded error on each individual cell would accumulate to order
$k$, too large for \eqref{eq:main-law}. We instead cut into geometric
blocks, use a uniform finite-chain Floquet estimate, and control
the total variation between block endpoints by summation by parts.
This proves \eqref{eq:intro-reduction} without assuming that the
periodic density of states is smooth at band edges.

Section~\ref{sec:forms} establishes the source form and cut comparison.
Section~\ref{sec:periodic} supplies the periodic estimates, including
the uniform finite-chain error. The reduction is proved in
Section~\ref{sec:reduction}; Section~\ref{sec:coefficient} evaluates
and analyzes its second coefficient. Section~\ref{sec:general}
proves the asymptotic-coupling and endpoint statements by form
comparison, with the domain inclusions made explicit.
